% !TeX root = ../sections/02_exercises.tex

\begin{exercise}
	R-11. $V$ を整数を成分とする $n$ 次列ベクトル全体がなす加法群とする。
	\begin{enumerate}
		\item $V$ は階数 $n$ の $\mathbb{Z}$-自由加群であることを示せ。
		\item $A$ を整数成分の $n$ 次正方行列とし，自己準同型写像
		\[
		f:V\longrightarrow V
		\]
		を
		\[
		f(v)=Av
		\qquad
		(v\in V)
		\]
		で定める。
		このとき，$f$ が同型となるための必要十分条件は
		\[
		\det A=\pm 1
		\]
		であることを示せ。
	\end{enumerate}
\end{exercise}

\begin{background}
	整数成分の $n$ 次列ベクトル全体は
	\[
	V=\mathbb{Z}^n
	\]
	である。
	標準基底を
	\[
	e_1,\dots,e_n
	\]
	とすると，任意の $v\in V$ は一意的に
	\[
	v=a_1e_1+\cdots+a_ne_n
	\qquad
	(a_1,\dots,a_n\in\mathbb{Z})
	\]
	と表される。
	したがって $V$ は階数 $n$ の $\mathbb{Z}$-自由加群である。
	
	整数行列 $A$ によって定まる写像
	\[
	f(v)=Av
	\]
	は $\mathbb{Z}$-加群準同型である。
	
	この写像が同型であるとは，逆写像も $\mathbb{Z}$-加群準同型として存在することをいう。
	つまり，逆行列 $A^{-1}$ が存在し，しかもその成分がすべて整数でなければならない。
\end{background}

\begin{strategy}
	まず，$V=\mathbb{Z}^n$ の標準基底
	\[
	e_1,\dots,e_n
	\]
	を用いれば，$V$ が階数 $n$ の自由 $\mathbb{Z}$-加群であることは直ちに分かる。
	
	次に，$f(v)=Av$ が同型であると仮定する。
	このとき，逆写像も整数行列によって表される。
	すなわち，ある整数行列 $B$ が存在して
	\[
	AB=BA=I_n
	\]
	となる。
	両辺の行列式を取ると
	\[
	\det A\cdot \det B=1
	\]
	である。
	ここで $\det A,\det B$ は整数なので，
	\[
	\det A=\pm 1
	\]
	でなければならない。
	
	逆に，
	\[
	\det A=\pm 1
	\]
	とする。
	余因子行列を用いると
	\[
	A^{-1}
	=
	\frac{1}{\det A}\operatorname{adj}(A)
	\]
	である。
	$\operatorname{adj}(A)$ は整数行列であり，$\det A=\pm 1$ だから，
	$A^{-1}$ も整数行列である。
	
	したがって，$A^{-1}$ によって逆写像が定まり，
	$f$ は $\mathbb{Z}$-加群の同型である。
\end{strategy}